\section{Mass-action parametrization and the exact decision formula}
\label{sec:semialgebraic-independent}

Consider indexed reactions $y_r\to y'_r$, with source matrix
$Y=[y_1\ \cdots\ y_m]$ and stoichiometric matrix
$\Gamma=[y'_1-y_1\ \cdots\ y'_m-y_m]$.  At a positive state $x^*$, put
\[
 v_r=k_r(x^*)^{y_r},\qquad h_i=(x_i^*)^{-1}.
\]
The equilibrium equation is $\Gamma v=0$.  Differentiating
$k_rx^{y_r}$ at $x^*$ gives the row vector
$v_ry_r^T\operatorname{diag}(h)$, and hence
\begin{equation}
 J(v,h)=\Gamma\operatorname{diag}(v)Y^T\operatorname{diag}(h).
 \label{eq:mass-action-factor-independent}
\end{equation}
Conversely, if $v>0$, $\Gamma v=0$, and $h>0$, choose
\[
 x_i^*=h_i^{-1},\qquad
 k_r=\frac{v_r}{(x^*)^{y_r}}.
\]
These quantities are positive, $x^*$ is an equilibrium, and direct
differentiation recovers \eqref{eq:mass-action-factor-independent}.  Thus the
parametrization is onto the set of positive-equilibrium Jacobians; it also
handles the zero complex, parallel indexed reactions, autocatalysis, and
arbitrary molecularity without modification.

Set
\[
 A(v)=\Gamma\operatorname{diag}(v)Y^T,
 \qquad J(v,h)=A(v)\operatorname{diag}(h).
\]
For $I\subseteq[n]$, column scaling gives
\begin{equation}
 (-1)^{|I|}\det J(v,h)[I,I]
 =\left(\prod_{i\in I}h_i\right)
   (-1)^{|I|}\det A(v)[I,I].
 \label{eq:minor-scale-independent}
\end{equation}

Let $B$ and $C$ be rational matrices with columns of $B$ forming a basis of
$S=\operatorname{im}\Gamma$ and $CB=I$.  Define
\[
 R(v,h)=CJ(v,h)B.
\]
By Section~\ref{sec:conservation-independent}, this is a coordinate matrix for
$J(v,h)|_S$ and changes only by similarity when the rational basis changes.

\begin{theorem}[Exact all-mobile characterization]
\label{thm:semialgebraic-independent}
A finite classical mass-action network admits a positive equilibrium that is
Hurwitz on $S$ and a strictly positive all-species diffusion matrix producing
a positive real nonzero-mode eigenvalue if and only if there are
$v\in\mathbb R_{>0}^m$, $h\in\mathbb R_{>0}^n$, and a nonempty
$I\subseteq[n]$ such that
\begin{align}
 \Gamma v&=0,\label{eq:char-flux-independent}\\
 R(v,h)&\text{ is Hurwitz},\label{eq:char-hurwitz-independent}\\
 (-1)^{|I|}\det A(v)[I,I]&<0.
 \label{eq:char-minor-independent}
\end{align}
\end{theorem}

\begin{proof}
For necessity, begin with positive rates and equilibrium.  The construction
above gives positive $v,h$ satisfying \eqref{eq:char-flux-independent}, while
relative homogeneous stability is exactly
\eqref{eq:char-hurwitz-independent}.  The spatial positive real eigenvalue and
Theorem~\ref{thm:fixedJ-independent} give a negative signed principal minor of
$J$.  Equation~\eqref{eq:minor-scale-independent} removes $h$ from its sign and
gives \eqref{eq:char-minor-independent}.  Under relative Hurwitz stability the
empty and full sets cannot be witnesses, so a nonempty branch is sufficient.

For sufficiency, reconstruct $x^*$ and $k$ from $v,h$.  Condition
\eqref{eq:char-hurwitz-independent} gives homogeneous stability on the
compatibility class.  Equations \eqref{eq:minor-scale-independent} and
\eqref{eq:char-minor-independent} give a negative signed principal minor of the
full Jacobian.  Theorem~\ref{thm:fixedJ-independent} supplies $D\succ0$ and a
positive real eigenvalue of $J-D$.  Rescaling $D$ by any chosen positive
Neumann eigenvalue realizes it as $q^2D_{\rm phys}$ for a nonzero spatial mode.
The amplitude is unrestricted by conservation, as proved above.
\end{proof}

Routh--Hurwitz determinants of the characteristic polynomial of $R(v,h)$ turn
\eqref{eq:char-hurwitz-independent} into finitely many strict rational
polynomial inequalities.  Together with the linear equations and positivity,
Theorem~\ref{thm:semialgebraic-independent} is a finite semialgebraic formula
computed from the indexed pair $(\Gamma,Y)$.  It eliminates the original rate
constants, the equilibrium concentrations, the diffusion coefficients, the
domain scale, and the wave number, but it does not assert a first, simple, or
transverse stationary crossing.

For a designated mobile set $M$, replace
\eqref{eq:char-minor-independent} by the finite disjunction
\[
 \exists\lambda>0,\quad I\supseteq M^c:\quad
 \det(\lambda I_{|I|}-J(v,h)[I,I])<0.
\]
Proposition~\ref{prop:conservation-splitting-independent} verifies the
background positivity required by Theorem~\ref{thm:mobile-independent}.
